\begin{proposition}[A unique transverse family type and its stable tangent bundle]\hypertarget{tag-02MR}{}\label{maps:src:E.2.7}
\StatementRecord{maps:src:E.2.7}{proved}{02MR}{\cref{maps:src:E.2.6}}


For fixed marked sphere fibrations, all M\"obius transformations
whose critical triples are disjoint give orientation-preservingly
diffeomorphic shifted products. The same conclusion holds while
either sphere varies in its proper smooth marked period family.
For every such fivefold $P$ there is an isomorphism of real vector
bundles
\begin{equation}
 TP\oplus\mathbb R^4\simeq\mathbb R^{14}.
 \hypertarget{tag-02BI}{}\label{maps:eq:E.MV9}
\end{equation}
Thus all its positive-degree Stiefel--Whitney and Pontryagin
classes vanish. For its given complex tangent bundle,
$c_1(P)=c_2(P)=0$ integrally.
\end{proposition}


\begin{proof}
\begin{enumerate}[label=\arabic*.,leftmargin=*]
\item The allowed transformations form the complement in
$\operatorname{PGL}_2(\mathbb C)$ of the nine analytic
hypersurfaces imposing an equality of one critical value from
each triple. This complement is path connected: a piecewise
smooth path between two points can be perturbed relative to
its endpoints to avoid these sets, which have real codimension
two. Equivalently, use local complex coordinate balls and
perturb each segment there. Choose a smooth such path.

\item The corresponding fibre products form a compact smooth family
over its parameter interval. Indeed the vertical derivative of
$f_1(x)-f_2(y)$ is surjective at every point, since its two
factors cannot be critical at the same base value. The same
surjectivity solves the equation for a lift of the interval
coordinate vector field. A partition of unity supplies a
global smooth lift; compactness makes its flow exist for the
whole interval. This flow gives the claimed diffeomorphisms.
The complex orientations vary continuously and hence are
preserved. The identical argument applies to smooth proper
period families of either factor, keeping the critical triples
disjoint throughout.

\item View $P$ as the transverse inverse image of the diagonal in
$B\times B$ under $f_1\times f_2$. Its real oriented normal
two-plane bundle is the underlying bundle of $g^*T_B$,
where $g:P\to B$ is the common base map. It is trivial:
its first Chern class is
\[
 c_1(g^*T_B)=2g^*[B]^\vee
    =2\operatorname{pr}_1^*f_1^*[B]^\vee=0,
\]
because $H^2(X;\mathbb Z)=0$, and complex line bundles are
classified by their first Chern class. The tangent bundles
of the two six-spheres become trivial after adding one real
line each. Restriction of their product tangent bundle gives
\[
 (TP\oplus\mathbb R^2)\oplus\mathbb R^2
       \simeq\mathbb R^{14},
\]
proving \eqref{maps:eq:E.MV9} and the characteristic-class assertions for
the real tangent bundle. Finally the holomorphic tangent
sequence is
\[
 0\longrightarrow TP\longrightarrow
 \operatorname{pr}_1^*TX\oplus\operatorname{pr}_2^*TX'
 \longrightarrow g^*T_B\longrightarrow0.
\]
Both sphere factors have integral $c_1=c_2=0$, since their
cohomology in degrees two and four vanishes. The last line
has zero first Chern class, as just proved. Multiplying total
Chern classes in this exact sequence gives
$c_1(P)=c_2(P)=0$ integrally.
\end{enumerate}
\end{proof}